\chapter*{About This Guide}
\addcontentsline{toc}{chapter}{About This Guide}

This guide accompanies the \emph{DSU dsDNA Core uafR Undergraduate Research
Training Program: Student Manual}. It is written for instructors, research
mentors, and program staff who are responsible for preparing the learning
environment, reviewing student evidence, and deciding when a project is ready
to advance through the program gates.

\section*{What Is Shared Publicly}
The public curriculum includes the student manual, this generic instructor
guide, training scripts, worksheets, rubrics, simulated examples, and
installation tools. This transparency allows students to understand how their
work will be evaluated and allows other research programs to inspect the
curriculum before adopting it.

\section*{What Belongs in a Local Supplement}
The public guide does not define institution-specific employment rules,
compensation, schedules, accommodations, emergency procedures, unpublished
project details, credentials, or restricted-data locations. Before the program
begins, those items should be documented in a locally approved supplement and
reviewed with every participant.

\section*{Companion Documents}
\begin{itemize}
  \item The student manual contains the complete 70-day curriculum, labs,
        readings, project tracks, worksheets, rubrics, and student setup guide.
  \item The package README and pkgdown site contain current function-level
        documentation and species-first plant phytochemistry examples.
  \item The offline student bundle provides a fixed package archive, the student
        manual, scripts, integrity checks, and an acceptance test.
  \item Repository workflows verify that public manuals build and that the
        offline training path still runs against the current package.
\end{itemize}

\begin{goldbox}{Support Boundary}
The curriculum teaches evidence-aware analysis; it does not replace laboratory
safety training, responsible-conduct requirements, data-governance approval,
chemical identity confirmation, statistical consultation, or project-specific
scientific supervision.
\end{goldbox}
